\documentclass[10pt,a4paper]{amsart}
\usepackage[margin=19mm]{geometry}
\usepackage{amsmath,amssymb,mathtools}
\usepackage[scr=rsfs,cal=euler]{mathalfa}
\usepackage{xfrac}
\usepackage[unicode,hidelinks,bookmarks=false]{hyperref}
\newcommand{\ie}{i.e.\ }
\newcommand{\cf}{cf.\ }
\newcommand{\resp}{respectively\ }
\newcommand{\Subsection}{\subsection}
\providecommand{\nobreakdash}{\nobreak\hbox{-}}
\setlength{\emergencystretch}{2em}
\multlinegap0pt
\usepackage{bm}
\usepackage{bbm}
\usepackage{dsfont}
\usepackage{xspace}
\usepackage{cancel}
\usepackage{mathtools}
\usepackage{amsthm}
\makeatletter
\@ifundefined{thm}{\newtheorem{thm}{Theorem}}{}
\makeatother
\makeatletter
\expandafter\ifx\csname proofa\endcsname\relax
\newenvironment{proofa}{\begin{proof}[Proof of Theorem 1.3]}{\end{proof}}
\fi
\expandafter\ifx\csname proofb\endcsname\relax
\newenvironment{proofb}{\begin{proof}[Proof of Theorem 1.2]}{\end{proof}}
\fi
\makeatother
\title[The ordering of hypertrees and unicyclic hypergraphs by the ]{The ordering of hypertrees and unicyclic hypergraphs by the traces of $\mathcal{A}\_{\alpha}$-tensor}
\author{Jueru Liu, Lizhu Sun, Changjiang Bu}
\date{Source: arXiv:2507.02650v1 (CC BY 4.0)}
\begin{document}
\maketitle

\noindent\emph{Source and licence.} The ordering of hypertrees and unicyclic hypergraphs by the traces of $\mathcal{A}_{\alpha}$-tensor. Version arXiv:2507.02650v1 (CC BY 4.0), distributed under the Creative Commons Attribution 4.0 International licence, as recorded in the arXiv licence metadata for this exact version and shown on the abstract page https://arxiv.org/abs/2507.02650v1. Full rights notice: licenses/arxiv-2507.02650-CC-BY-4.0.txt.\par\smallskip

\noindent\textit{Editorial scope.} Counted unit: the Thm whose source label is \texttt{None} in the article (source0.tex lines 603--605, source label \texttt{None}), delivered together with the article's own complete original proof of it (source0.tex lines 607--652, one \texttt{proof} environment of 46 lines). No mathematics is written, repaired or re-derived: statement and proof are the article's own text line for line. Required context retained uncounted: T (lines 241--250); cgpn-g (lines 575--577); firstcycle (lines 588--590). Every definition, display and label of those lines is retained unchanged. Omitted from this package and not redistributed: 1--240, 251--574, 578--587, 591--602, 606--606, 653--805. Nothing inside a retained excerpt is omitted or altered: every retained body is delivered exactly as the article's lines are written, so no omission receipt of any kind accompanies this package. Citations: the retained excerpts carry no citation command at all, so no bibliography entry of the article is transcribed and no reference is redistributed. Reference adaptation: none was needed and none was made; every reference inside the delivered unit resolves inside it, so no unresolved reference or citation is produced. Printed slips kept literal: no printed word, symbol, hypothesis or number of the counted statement or of its proof was altered. Layout and class: the article's own class is \texttt{elsarticle1} with packages including subfig, graphicx, caption, color, amssymb, amsthm, amsmath, epic, setspace, float, multirow, hyperref; the wrapper is the lane's pinned \texttt{amsart} editorial wrapper at 10pt on A4, which restates the theorem-like environments the retained text needs and adds the article's own macro definitions that the retained text needs (none), with extra wrapper packages (bm, bbm, dsfont, xspace, cancel, mathtools). The delivered numbering is the unit's own. Assets: no retained span contains a figure environment, a graphics-inclusion command, a PDF-page inclusion, a table or a TikZ picture; no image file is copied into this package, the package declares no asset dependency and no image file is redistributed with it. Licence: version arXiv:2507.02650v1 of this article is distributed under the Creative Commons Attribution 4.0 International licence, as recorded in the arXiv licence metadata for this exact version and shown on the abstract page https://arxiv.org/abs/2507.02650v1. A full rights notice is delivered at licenses/arxiv-2507.02650-CC-BY-4.0.txt. Characters: every retained excerpt is pure ASCII, so the delivered engine, XeLaTeX with its default Latin Modern text font, reports no missing character.
\begin{thm}\label{T}
Let $k\ge2$ and $\mathcal{H}$ be a $k$-uniform hypergraph with degree sequence $d_1,d_2,\ldots,d_n$.
Then
\begin{equation*}
\begin{aligned}
&\mathrm{Tr}_{d}(\mathcal{A}_{\alpha}(\mathcal{H}))=\\& \phi(d)+(1-\alpha)^{d}\mathrm{Tr}_{d}(\mathcal{A}(\mathcal{H}))+d(k-1)^{n}\sum\limits_{H\in\mathcal{V}_{d-1}(\mathcal{H})}\sum\limits_{f\in\mathcal{F}_{d}(H)}\frac{\tau(f)\pi_{f}(\mathcal{A}_{\alpha}(\mathcal{H}))}{\prod\limits_{v\in V(f)}d^{+}(v)},
\end{aligned}
\end{equation*}
where $\tau(f)$ is the number of rooted spanning trees of the multi-digraph $D(f)$ with a specified root.
\end{thm}

\begin{thm}\label{cgpn-g}
Among all $k$-uniform linear unicyclic hypergraphs with $m\ge3$ hyperedges and girth $g(3\le g\le m)$, the first hypergraph in $S_{\alpha}$-order is $C_{g}^{(k)}\cdot P_{m-g}^{(k)}$ for $k\ge3$ and $0<\alpha<1$.
\end{thm}

\begin{thm}\label{firstcycle}
Among all $k$-uniform linear unicyclic hypergraphs with $m\ge3$ hyperedges, the first hypergraph in $S_{\alpha}$-order is the hypercycle $C_{m}^{(k)}$ for $k\ge3$ and $0<\alpha<1$.
\end{thm}

\begin{thm}
Among all $k$-uniform linear unicyclic hypergraphs with $m\ge3$ hyperedges, the second hypergraph in $S_{\alpha}$-order is $C_{m-1}^{(k)}\cdot P_{1}^{(k)}$ for $k\ge3$ and $0<\alpha<1$.
\end{thm}

\begin{proof}
From Theorem \ref{cgpn-g} and Theorem \ref{firstcycle}, we know that the second hypergraph in $S_{\alpha}$-order among all $k$-uniform linear unicyclic hypergraphs with $m$ hyperedges is the first in $\mathbb{U}_{m}\setminus\{C_{m}^{(k)}\}$.
For any $U_1, U_2\in\mathbb{U}_{m}\setminus\{C_{m}^{(k)}\}$ and $d=1,\ldots,2k+1$, we can know that $\mathcal{V}_{d}(U_1)=\mathcal{V}_{d}(U_2)$ and $\pi_{f}(\mathcal{A}_{\alpha}(U_1))=\pi_{f}(\mathcal{A}_{\alpha}(U_2))$ for any $f\in\mathcal{F}_{d}(H)$ and $H\in\mathcal{V}_{d-1}(U_1)$.
The first $2k+1$ orders spectral moments and the degree sequences of all hypergraphs in $\mathbb{U}_{m}\setminus\{C_{m}^{(k)}\}$ are same.
From Theorem \ref{T}, we know that $\mathrm{Tr}_{d}(\mathcal{A}_{\alpha}(U_1))=\mathrm{Tr}_{d}(\mathcal{A}_{\alpha}(U_2))$ for any $U_1, U_2\in\mathbb{U}_{m}\setminus\{C_{m}^{(k)}\}$ and $d\in[2k+1]$.

Next, we compare the $2k+2$-nd order $\mathcal{A}_{\alpha}$-spectral moments of hypergraphs in $\mathbb{U}_{m}\setminus\{C_{m}^{(k)}\}$.
For any $U\in\mathbb{U}_{m}\setminus\{C_{m}^{(k)},C_{m-1}^{(k)}\cdot P_{1}^{(k)}\}$, we can see that $\mathrm{Tr}_{2k+2}(\mathcal{A}(U))=\mathrm{Tr}_{2k+2}(\mathcal{A}(C_{m-1}^{(k)}\cdot P_{1}^{(k)}))$.
Since $U$ and $C_{m-1}^{(k)}\cdot P_{1}^{(k)}$ have the same degree sequence and the differences between $\pi_{f}(\mathcal{A}_{\alpha}(U))$ and $\pi_{f}(\mathcal{A}_{\alpha}(C_{m-1}^{(k)}\cdot P_{1}^{(k)}))$ are the terms obtained from $f\in\mathcal{F}_{2k+2}(P_2^{(k)})$, from Theorem \ref{T}, we have
\begin{align*}
&\mathrm{Tr}_{2k+2}(\mathcal{A}_{\alpha}(C_{m-1}^{(k)}\cdot P_{1}^{(k)}))-\mathrm{Tr}_{2k+2}(\mathcal{A}_{\alpha}(U))\\
=&(2k+2)(k-1)^{m(k-1)}\sum\limits_{f\in\mathcal{F}_{2k+2}(P_2^{(k)})}\left( \frac{\tau(f)\pi_{f}(\mathcal{A}_{\alpha}(C_{m-1}^{(k)}\cdot P_{1}^{(k)}))}{\prod\limits_{v\in V(f)}d^{+}(v)} - \frac{\tau(f)\pi_{f}(\mathcal{A}_{\alpha}(U))}{\prod\limits_{v\in V(f)}d^{+}(v)} \right).
\end{align*}

Given a $k$-uniform hypergraph $\mathcal{H}$, for $P_{2}^{(k)}\subset \mathcal{H}$ with the vertex set $V(P_{2}^{(k)})=\{i_1,\ldots,i_{2k-1}\}$ and the hyperedge set $E(P_2^{(k)})=\{ e_1, e_2 \}$, where $i_1<\cdots<i_{2k-1}$, $e_1=\{i_1,\ldots,i_k\}$ and $e_2=\{i_k,\ldots,i_{2k-1}\}$.
Let $\alpha_j$ be a permutation of $e_1 \setminus\{i_j\}$ for $j=1,\ldots,k-1$ and $\alpha_l$ be a permutation of $e_2 \setminus\{i_l\}$ for $l=k+1,\ldots,2k-1$.
Let $\alpha_k$ be a permutation of $e_1 \setminus\{i_k\}$ and $\alpha_{k+1}$ be a permutation of $e_2 \setminus\{i_k\}$.
Let $\beta_j=i_j\cdots i_j\in[N]^{k-1}$ for $j=1,\ldots,2k-1$, where $N=m(k-1)$.
The elements in $\mathcal{F}_{2k+2}(P_{2}^{(k)})$ have the following four cases.\\
\textbf{Case 1.} Let $f_1=(i_1\alpha_1,\ldots,i_k\alpha_k,i_k\beta_k,i_k\beta_k,i_k\alpha_{k+1},\ldots,i_{2k-1}\alpha_{2k-1})$.
Then $$\pi_{f_1}(\mathcal{A}_{\alpha}(\mathcal{H}))=\alpha^2(1-\alpha)^{2k}\left( \frac{1}{(k-1)!} \right)^{2k}d_{i_k}^{2},~ \sum_{v\in V(f_1)}d_v^{+}=4(k-1)^{2k-1}.$$
And there are $12((k-1)!)^{2k}$ elements in $\mathcal{F}_{2k+2}(P_{2}^{(k)})$ that have the same induced multi-digraph as $f_1$.\\
\textbf{Case 2.} Let $f_2=(i_1\alpha_1,\ldots,i_j\alpha_j,i_j\beta_j,i_j\beta_j,i_{j+1}\alpha_{j+1},\ldots,i_k\alpha_k,i_k\alpha_{k+1},\ldots,i_{2k-1}\alpha_{2k-1})$, where $j\in[2k-1]\setminus\{k\}$. Then $$\pi_{f_2}(\mathcal{A}_{\alpha}(\mathcal{H}))=\alpha^2(1-\alpha)^{2k}\left( \frac{1}{(k-1)!} \right)^{2k}d_{i_j}^{2},~ \sum_{v\in V(f_2)}d_v^{+}=6(k-1)^{2k-1}.$$
And there are $6((k-1)!)^{2k}$ elements in $\mathcal{F}_{2k+2}(P_{2}^{(k)})$ that have the same induced multi-digraph as $f_2$.\\
\textbf{Case 3.} Let $f_3=(i_1\alpha_1,\ldots,i_j\alpha_j,i_j\beta_j,i_{j+1}\alpha_{j+1},\ldots,i_k\alpha_k,i_k\beta_k,i_k\alpha_{k+1},\ldots,i_{2k-1}\alpha_{2k-1})$, where $j\in[2k-1]\setminus\{k\}$. Then $$\pi_{f_3}(\mathcal{A}_{\alpha}(\mathcal{H}))=\alpha^2(1-\alpha)^{2k}\left( \frac{1}{(k-1)!} \right)^{2k}d_{i_j}d_{i_k},~ \sum_{v\in V(f_3)}d_v^{+}=6(k-1)^{2k-1}.$$
And there are $12((k-1)!)^{2k}$ elements in $\mathcal{F}_{2k+2}(P_{2}^{(k)})$ that have the same induced multi-digraph as $f_3$.\\
\textbf{Case 4.} Let $f_4=(i_1\alpha_1,\ldots,i_j\alpha_j,i_j\beta_j,i_{j+1}\alpha_{j+1},\ldots,i_l\alpha_l,i_l\beta_l,i_{l+1}\alpha_{l+1},\ldots,i_k\alpha_k, i_k\alpha_{k+1},\\ \ldots,i_{2k-1}\alpha_{2k-1})$, where $j\ne l\in[2k-1]\setminus\{k\}$. Then $$\pi_{f_4}(\mathcal{A}_{\alpha}(\mathcal{H}))=\alpha^2(1-\alpha)^{2k}\left( \frac{1}{(k-1)!} \right)^{2k}d_{i_j}d_{i_l},~ \sum_{v\in V(f_4)}d_v^{+}=8(k-1)^{2k-1}.$$
And there are $8((k-1)!)^{2k}$ elements in $\mathcal{F}_{2k+2}(P_{2}^{(k)})$ that have the same induced multi-digraph as $f_4$.

It can be seen that the digraph obtained by removing all loops of $D_{f_i}$ is the digraph obtained by intersecting two complete digraphs with $k$ vertices at a common vertex, then $\tau(f_i)=k^{2k-4}$ for $i\in[4]$.

For any $U\in\mathbb{U}_{m}\setminus\{C_{m}^{(k)},C_{m-1}^{(k)}\cdot P_{1}^{(k)}\}$, by the vertex-hyperedge incidence relations of hypergraphs $C_{m-1}^{(k)}\cdot P_{1}^{(k)}$ and $U$, from the above statements,
when $0<\alpha<1$, we have
\begin{align*}
&\sum\limits_{f\in\mathcal{F}_{2k+2}(P_2^{(k)})}\left( \frac{\tau(f)\pi_{f}(\mathcal{A}_{\alpha}(C_{m-1}^{(k)}\cdot P_{1}^{(k)}))}{\prod\limits_{v\in V(f)}d^{+}(v)} - \frac{\tau(f)\pi_{f}(\mathcal{A}_{\alpha}(U))}{\prod\limits_{v\in V(f)}d^{+}(v)} \right).\\
=&(k-1)^{1-2k}(k-1)^{2k-4} \alpha^{2}(1-\alpha)^{2k} \Bigg( \sum\limits_{P_{2}^{k}\subset C_{m-1}^{(k)}\cdot P_{1}^{(k)}}\bigg( \sum\limits_{i\in V(P_{2}^{(k)})}d_i^2+\sum\limits_{i,j\in V(P_{2}^{(k)})}d_i d_j \bigg)\\
&-\sum\limits_{P_{2}^{k}\subset U}\bigg( \sum\limits_{i\in V(P_{2}^{(k)})}d_i^2+\sum\limits_{i,j\in V(P_{2}^{(k)})}d_i d_j \bigg) \Bigg)\\
=&(k-1)^{1-2k}(k-1)^{2k-4} \alpha^{2}(1-\alpha)^{2k} \Bigg( \left( 2\binom{2k-4}{2}+14(2k-4)+48 \right)\\
&-\left( \binom{2k-5}{2}+\binom{2k-3}{2}+9(2k-5)+5(2k-3)+52 \right) \Bigg)\\
=&-(k-1)^{1-2k}(k-1)^{2k-4}\alpha^{2}(1-\alpha)^{2k}<0.
\end{align*}
Thus, $\mathrm{Tr}_{2k+2}(\mathcal{A}_{\alpha}(C_{m-1}^{(k)}\cdot P_{1}^{(k)}))<\mathrm{Tr}_{2k+2}(\mathcal{A}_{\alpha}(U))$ for any $U\in\mathbb{U}_{m}\setminus\{C_{m}^{(k)}\}$ with $U\neq C_{m-1}^{(k)}\cdot P_{1}^{(k)}$.
It follows that $C_{m-1}^{(k)}\cdot P_{1}^{(k)}\preceq_{\alpha}U$ for all $U\in\mathbb{U}_{m}\setminus\{C_{m}^{(k)}\}$ with the equality if and only if $U\cong C_{m-1}^{(k)}\cdot P_{1}^{(k)}$.
Then $C_{m-1}^{(k)}\cdot P_{1}^{(k)}$ is the first hypergraph in $\mathbb{U}_{m}\setminus\{C_{m}^{(k)}\}$.
Therefore, $C_{m-1}^{(k)}\cdot P_{1}^{(k)}$ is the second hypergraph in $S_{\alpha}$-order among all $k$-uniform linear unicyclic hypergraphs with $m$ hyperedges.
\end{proof}

\end{document}
